\documentclass[10pt]{article}
\usepackage[margin=1in]{geometry}
\usepackage[T1]{fontenc}
\usepackage[utf8]{inputenc}
\usepackage{lmodern}
\usepackage{amsmath,amssymb,amsthm}
\usepackage{booktabs}
\usepackage{pgfplots}
\pgfplotsset{compat=1.18}
\usepackage{xcolor}
\usepackage[hidelinks]{hyperref}
\usepackage{microtype}

\definecolor{QRBlue}{HTML}{1462A6}
\definecolor{SchurOrange}{HTML}{CB551D}
\definecolor{Muted}{HTML}{586572}

\input{data/numbers_v2.tex}

\newcommand{\norm}[1]{\left\lVert #1\right\rVert}
\newcommand{\kap}{\kappa_2}
\newcommand{\plotqr}[1]{%
  \addplot+[color=QRBlue,mark=*,mark size=1.7pt,line width=.9pt,
  error bars/y dir=both,error bars/y explicit,error bars/error bar style={line width=.5pt}]
  table[x=x,y=y,y error plus=plus,y error minus=minus,col sep=comma]{data/#1};}
\newcommand{\plotschur}[1]{%
  \addplot+[color=SchurOrange,mark=square*,mark size=1.6pt,line width=.9pt,
  error bars/y dir=both,error bars/y explicit,error bars/error bar style={line width=.5pt}]
  table[x=x,y=y,y error plus=plus,y error minus=minus,col sep=comma]{data/#1};}
\pgfplotsset{every axis/.append style={
  axis lines=left,tick align=outside,
  label style={font=\small},tick label style={font=\footnotesize},
  title style={font=\small\bfseries},
  legend style={font=\footnotesize,draw=none,fill=white,cells={anchor=west}},
  grid=major,grid style={gray!20},major tick length=2pt,
  scaled ticks=false,unbounded coords=jump}}


\title{Dense Direct Solvers for Equality Constrained\\Quadratic Optimization\\
\large A QR Null Space Method versus a Schur Complement Method}
\author{Efe Cinar\\CPSC 542G, Midterm Project}
\date{November 2, 2026}

\begin{document}
\maketitle

\begin{abstract}
I compare two dense direct methods for the equality constrained quadratic program
$\min \tfrac12 x^\top Hx - g^\top x$ subject to $Cx=d$ with $H$ symmetric positive definite.
The first is a null space method built on a full QR factorization of $C^\top$, the second
a Schur complement method built on two Cholesky factorizations. Both solve the same KKT system in exact
arithmetic and both cost $O(n^3)$ at fixed constraint fraction $p/n$, but they distribute
the work and the rounding error differently. On dense problems with up
to $n=12000$ unknowns, on three problems taken from applications (a CUTEst optimal control
problem, a SuiteSparse power network least squares problem, and ordinary kriging on real
elevation data), and on a conditioning study with certified high precision references,
I find that the Schur method is consistently faster, by a factor of about~2 at
$p/n=0.5$ and 6 to 7 when $p\ll n$, with no advantage left at $p/n=0.9$. The QR method is
consistently more accurate, by one order of magnitude on well conditioned problems and
seven orders at $\kap(H)=10^8$, and it is the only one of the two that survives
$\kap(C)\ge 10^{10}$, where the Schur complement loses positive definiteness in double
precision. A multiprecision rerun of the failed cases confirms that this is a
finite precision effect of forming $S=CH^{-1}C^\top$ explicitly, not a property of the
problems.
\end{abstract}

\section{Introduction}

Equality-constrained quadratic programs (EQPs) are the basic building block of
constrained optimization. Every step of a sequential quadratic programming or
interior point method solves one, and the equality constrained case is also the simplest
setting in which the interaction between an objective and a set of constraints can be
studied with linear algebra alone. The problem considered here is
\begin{equation}\label{eq:qp}
\min_{x\in\mathbb R^n}\ \tfrac12 x^\top H x - g^\top x
\qquad\text{subject to}\qquad Cx=d,
\end{equation}
with $H=H^\top\succ 0$ and $C\in\mathbb R^{p\times n}$ of full row rank, $p<n$.
Under these assumptions the problem has a unique solution, and the two classical dense
approaches to computing it are to eliminate either the unknown $x$ (the Schur complement,
or range-space, method) or the constraint (the null space method). Both are described in
the survey of Benzi, Golub and Liesen~\cite{benzi2005}.

The goal of this project was to implement both methods from scratch, measure what they
actually cost on dense matrices of realistic size, and find out how they behave when $H$
or $C$ is ill conditioned. The two methods are often presented as interchangeable, and the
operation counts alone do not say which one to prefer. I was curious whether the
numerical behaviour would separate them. It does, and in a way that the operation counts
do not predict.

The remainder of the report derives the two methods and their cost (Section~\ref{sec:methods}),
describes the test problems, references and timing protocol (Section~\ref{sec:setup}),
presents the results (Section~\ref{sec:results}) and discusses what they do and do not
establish (Section~\ref{sec:discussion}). The code is available on GitHub, see
Section~\ref{sec:code}.

\section{The two methods}\label{sec:methods}

\subsection{KKT conditions}

With the Lagrangian $\mathcal L(x,\lambda)=\tfrac12 x^\top Hx - g^\top x+\lambda^\top(Cx-d)$,
the first order conditions $\nabla_x\mathcal L=0$, $\nabla_\lambda\mathcal L=0$ are the linear
system
\begin{equation}\label{eq:kkt}
\begin{bmatrix} H & C^\top\\ C & 0\end{bmatrix}
\begin{bmatrix} x\\ \lambda\end{bmatrix}
=\begin{bmatrix} g\\ d\end{bmatrix},
\qquad K v = b .
\end{equation}
Because $H\succ0$ the objective is strictly convex on the affine set $\{Cx=d\}$, so the
minimizer is unique, and because $\operatorname{rank}C=p$ the multiplier $\lambda$ is
unique as well. The KKT matrix $K$ is symmetric but indefinite (its zero block forces $p$ negative
eigenvalues), so it cannot be factored by Cholesky directly. Both methods below reduce
\eqref{eq:kkt} to a smaller symmetric positive definite (SPD) system that can.

\subsection{Schur complement method}

Solving the first block row of \eqref{eq:kkt} for $x$ gives $x=H^{-1}(g-C^\top\lambda)$.
Substituting into the second row,
\begin{equation}\label{eq:schur}
\underbrace{CH^{-1}C^\top}_{S}\,\lambda = CH^{-1}g - d .
\end{equation}
The inverse is never formed. With the Cholesky factorization $H=L_HL_H^\top$, the
implementation computes $u=H^{-1}g$ and $W=H^{-1}C^\top$ by triangular solves (one
right hand side for $u$, $p$ for $W$), forms $S=CW$ and the right hand side $Cu-d$,
factors $S=L_SL_S^\top$, solves for $\lambda$, and recovers $x=u-W\lambda$.

$S$ is SPD. It is symmetric because $H^{-1}$ is, and for $\lambda\neq0$ we have
$\lambda^\top S\lambda=(C^\top\lambda)^\top H^{-1}(C^\top\lambda)>0$, since $C^\top$ has
full column rank and $H^{-1}\succ0$. Writing $H^{-1}=L_H^{-\top}L_H^{-1}$ also shows
\begin{equation}\label{eq:gram}
S = BB^\top,\qquad B = CL_H^{-\top},\qquad \kap(S)=\kap(B)^2 ,
\end{equation}
so $S$ is a Gram matrix whose condition number is the square of that of $B$. This is the
observation that will explain the failures in Section~\ref{sec:cond}.

\subsection{QR null space method}

The null space method keeps the constraint satisfied by construction. A full QR
factorization of the tall matrix $C^\top$,
\begin{equation}
C^\top=\begin{bmatrix}Q_1 & Z\end{bmatrix}\begin{bmatrix}R\\0\end{bmatrix}=Q_1R,
\qquad Q=\begin{bmatrix}Q_1 & Z\end{bmatrix}\ \text{orthogonal},
\end{equation}
provides an orthonormal basis $Q_1\in\mathbb R^{n\times p}$ of the row space of $C$ and an
orthonormal basis $Z\in\mathbb R^{n\times(n-p)}$ of its null space, $CZ=0$. A particular
feasible point is obtained from one triangular solve, $R^\top y=d$, $x_0=Q_1y$, since
$Cx_0=R^\top Q_1^\top Q_1y=d$. Every feasible point is then $x=x_0+Zz$ with
$z\in\mathbb R^{n-p}$ free, and substituting into the objective gives an unconstrained
quadratic in $z$ whose gradient vanishes when
\begin{equation}\label{eq:reduced}
\underbrace{Z^\top HZ}_{G}\,z = Z^\top(g-Hx_0).
\end{equation}
$G$ is SPD by the same argument as $S$ (for $z\ne0$, $z^\top Gz=(Zz)^\top H(Zz)>0$ because
$Z$ has full column rank), so it is factored by Cholesky. Finally $x=x_0+Zz$, and the
multiplier is recovered from the first KKT equation, $C^\top\lambda=g-Hx$, i.e.\
$Q_1R\lambda=g-Hx$, which after multiplication by $Q_1^\top$ is the upper triangular
solve $R\lambda=Q_1^\top(g-Hx)$.

Recovering $\lambda$ is not optional. The residual of \eqref{eq:kkt} can only be
evaluated when both $x$ and $\lambda$ are available, and the stationarity equation is
what certifies that $x$ is the minimizer rather than merely a feasible point.

\subsection{Cost}

Table~\ref{tab:cost} lists the dominant operations. With $q=n-p$, every row is cubic in
$n$ once $p$ is a fixed fraction of $n$, but the constants differ. At $p=n/2$ the Schur
method costs roughly $n^3/3$ (Cholesky of $H$) $+\,n^3$ (the $p$ solves) $+\,n^3/2$
(forming $CW$) $+\,n^3/24\approx1.9n^3$ flops, while the QR method costs roughly
$1.6n^3$ (Householder QR of $C^\top$ including the explicit formation of the full
$n\times n$ factor $Q$) $+\,n^3$ (the product $HZ$) $+\,n^3/2$ ($Z^\top(HZ)$)
$+\,n^3/24\approx3.1n^3$. Two things follow. First, Cholesky is the cheapest cubic
operation, and the Schur method is built almost entirely out of it, so it should be the
faster of the two at any fixed ratio. Second, the two methods depend on $p$ in opposite
ways. The Schur terms grow with $p$, the QR reduced system shrinks with $p$, but the
formation of $Q$ and of $HZ$ does not go away, so there is a floor under the QR cost.
Both expectations are tested in Section~\ref{sec:results}.

\begin{table}[t]
\centering\small
\begin{tabular}{@{}lll@{}}
\toprule
Method & Operation & Work\\
\midrule
Schur & Cholesky of $H$ & $n^3/3+O(n^2)$\\
      & Solve $HW=C^\top$, form $CW$ & $O(n^2p+np^2)$\\
      & Cholesky of $S$ & $p^3/3+O(p^2)$\\
\midrule
QR/null space & Full $Q$, including basis construction & $O(n^2p)$\\
      & Form $Z^\top HZ$ & $O(n^2q+nq^2)$\\
      & Cholesky of $G$ & $q^3/3+O(q^2)$\\
\bottomrule
\end{tabular}
\caption{Dominant dense operations of the two implementations, $q=n-p$.}
\label{tab:cost}
\end{table}

\subsection{Implementation}

Both solvers are short MATLAB functions built from \texttt{chol}, \texttt{qr} and
backslash on triangular factors. They share one input check (real, finite, dense,
symmetric $H$, consistent dimensions, $0<p<n$) and deliberately contain no safety net.
There is no explicit inverse, no SVD based null space, no regularization and no
fallback to a general KKT solver. If a Cholesky factorization fails the solver raises an
error, and the experiment records the failure as a result. The reduced matrices $G$ and $S$
are symmetrized as $(G+G^\top)/2$ before factorization. This removes the $10^{-16}$ level
asymmetry that matrix products introduce and is not a regularization. I verified both
solvers against an independent Python reimplementation (SciPy) on the stored test data,
with relative differences below $10^{-12}$.

\section{Experimental design}\label{sec:setup}

\subsection{Problems}

Three kinds of input were used, with different purposes.

\subsubsection*{Synthetic dense problems} $H=U\Lambda U^\top$ with $U$ a random orthogonal
matrix (from the QR of a Gaussian matrix) and $\Lambda$ log spaced between $1$ and
$1/\kap(H)$, and $C=W\Sigma V^\top$ with $W$, $V$ random orthogonal and $\Sigma$ log spaced so
that $\kap(C)$ is prescribed. A target pair $(x^\star,\lambda^\star)$ is drawn first and
$g=Hx^\star+C^\top\lambda^\star$, $d=Cx^\star$ are built from it, so the solution is known by
construction. These problems are used for timing (with $\kap(H)=100$, $\kap(C)=10$) and
for the ill conditioned size sweep (with $\kap(H)=10^8$, $\kap(C)=10^6$).

\subsubsection*{Application problems} To avoid relying on random matrices only, I also
used three problems with a model behind them. \emph{DTOC3} from the CUTEst collection
\cite{cutest,dtoc3} is a discrete time optimal control problem with two states and one
control per time step. The dynamics are linear equality constraints and the objective is
a sum of squares, so the Hessian is diagonal. Eliminating the fixed initial states gives
an EQP with $n=12000$, $p=8000$. Because the diagonal Hessian could in principle be
exploited by the library routines even in full storage, I also solved an \emph{equivalent
dense representation}. For a random orthogonal $U$, the change of variables $x=Uy$ gives
$\tilde H=U^\top HU$, $\tilde C=CU$, $\tilde g=U^\top g$, with $\tilde H$ and $\tilde C$
fully dense and the optimizer $y^\star=U^\top x^\star$. This is the same problem in rotated
coordinates (verified numerically to $10^{-15}$), not a new one, and it was solved for
five different $U$. \emph{GEMAT1} is a power network matrix from the SuiteSparse
collection \cite{suitesparse}. Following the construction of Scott and T\r{u}ma
\cite{scotttuma2022} for equality constrained least squares, the densest rows of the
(transposed, column scaled) matrix are taken as equality constraints and the remaining
rows form a least squares objective, which in the form \eqref{eq:qp} is $H=A^\top A$,
$g=A^\top b$, with $n=4929$ and $p\in\{5,20\}$. \emph{Kriging} is ordinary kriging of
elevation on the UCI 3D Road Network data set \cite{uci_road} (North Jutland, 397\,978
distinct locations). For $n$ random training sites and one target site, $H_{ij}=\exp(-\|s_i-s_j\|^2/2\ell^2)+\tau^2\delta_{ij}$
with $\ell=10$~km and nugget $\tau^2=10^{-3}$, $g_i$ is the covariance between site $i$ and
the target, and the single constraint $\mathbf 1^\top x=1$ makes the predictor unbiased.
$H$ is naturally dense and SPD, and $p=1$ is the extreme case of few constraints.

\subsubsection*{Conditioning study} Synthetic problems with $n=80$, $p=20$ and five
seeds, with $\kap(H)\in\{1,10^4,10^8,10^{10},10^{12}\}$ at $\kap(C)=10$ and
$\kap(C)\in\{1,10^4,10^6,10^8,10^{10},10^{12}\}$ at $\kap(H)=100$. The size is small so
that each stored problem can be solved to 96 decimal digits with \texttt{mpmath}
(checked against a 64-digit solve). These are the only certified references in the
study.

\subsection{References and error measures}

For every run I record the relative solution error and a scaled KKT residual,
\begin{equation}
e_x=\frac{\norm{\hat x-x_{\rm ref}}_2}{\norm{x_{\rm ref}}_2},\qquad
e_\lambda=\frac{\norm{\hat\lambda-\lambda_{\rm ref}}_2}{\norm{\lambda_{\rm ref}}_2},\qquad
\eta_{K,F}=\frac{\norm{K\hat v-b}_2}{\norm K_F\norm{\hat v}_2+\norm b_2},
\end{equation}
with $\norm K_F^2=\norm H_F^2+2\norm C_F^2$. What $x_{\rm ref}$ is differs by family and
matters for interpretation. For the synthetic problems it is the constructed target.
Since $g$ and $d$ are formed in double precision, the target is the exact solution of the
stored problem only up to roughly $\kap\,\varepsilon$, which is negligible for the
well conditioned timing families and significant for the ill conditioned ones. For DTOC3
it is a sparse double precision KKT solve, and for GEMAT1 an augmented least squares
solve in double, whose own accuracy is limited to about $10^{-8}$ because
$\kap(A^\top A)\approx10^9$. Kriging has no independent reference, so there I report the
agreement between the two methods and the residuals. Only the conditioning study has references accurate enough to
call $e_x$ a forward error.

\subsection{Timing protocol}

All experiments ran in MATLAB Online R2026b (Linux, Intel oneAPI MKL 2025.3). A timed
solve is a complete call of one solver, input checks and all factorizations included.
Matrix generation, loading, the orthogonal transformation, reference computation and
error evaluation are outside the timer. Each case gets one warm up call per method, then
seven timing batches in which the solver order alternates. A batch repeats the solve
until 0.2~s have elapsed (at most 100 times) and reports the mean, which only matters for
the smallest problems. Per case I report the median over batches, and across the five
random seeds the median and the min and max range. Repeated batches are repeated
measurements of one problem, not independent problems. The largest size, $n=12000$, is
the largest that fits the memory of the MATLAB Online session, where each $n\times n$
double matrix is 1.07~GiB.

\section{Results}\label{sec:results}

\subsection{Scaling with $n$, and the effect of conditioning on time}

\begin{figure}[t]
\centering
\begin{tikzpicture}
\begin{axis}[width=8.4cm,height=6cm,xmode=log,ymode=log,
 xlabel={$n$},ylabel={Complete solve time (s)},
 xtick={1000,2000,4000,8000,12000},xticklabels={1000,2000,4000,8000,12000},
 ytick={.01,.1,1,10},xmin=850,xmax=14500,ymin=.004,ymax=30,legend pos=south east]
\plotqr{v2_size_qr.csv}\plotschur{v2_size_schur.csv}
\addplot[domain=1000:12000,samples=2,color=Muted,dashed,line width=.6pt]{0.0217*(x/1000)^3};
\addplot+[color=QRBlue,only marks,mark=o,mark size=2.4pt,line width=.8pt]
 table[x=x,y=y,col sep=comma]{data/v2_illcond_qr.csv};
\addplot+[color=SchurOrange,only marks,mark=square,mark size=2.3pt,line width=.8pt]
 table[x=x,y=y,col sep=comma]{data/v2_illcond_schur.csv};
\legend{QR/null space,Schur,$n^3$ reference,{QR, $\kap(H)=10^8$},{Schur, $\kap(H)=10^8$}}
\end{axis}\end{tikzpicture}
\hfill
\begin{minipage}[b]{6.4cm}
\small
\begin{tabular}{@{}lcc@{}}
\toprule
$n=12000$, $p/n=0.5$ & \multicolumn{2}{c}{$\kap(H),\ \kap(C)$}\\
 & $10^2,\ 10$ & $10^8,\ 10^6$\\
\midrule
QR time & \SizeQRTop\ s & \IllQRTTop\ s\\
Schur time & \SizeSchurTop\ s & \IllSchurTTop\ s\\
QR $e_x$ & $\WellQRXTop$ & $\IllQRXTop$\\
Schur $e_x$ & $\WellSchurXTop$ & $\IllSchurXTop$\\
Schur $e_\lambda$ & $10^{-13}$ & \IllSchurLTop\%\\
QR $\eta_{K,F}$ & $3\times10^{-17}$ & $\IllQRKKTTop$\\
Schur $\eta_{K,F}$ & $3\times10^{-16}$ & $\IllSchurKKTTop$\\
\bottomrule
\end{tabular}
\vspace{3mm}
\end{minipage}
\caption{Left, solve time against $n$ at $p/n=0.5$. Filled markers have $\kap(H)=100$ and
$\kap(C)=10$ and show the median over five seeds and seven batches, with bars for the seed
min and max. Hollow markers repeat the same sizes with $\kap(H)=10^8$ and $\kap(C)=10^6$ (two
seeds, one solve each). Right, times and errors at the largest size for both conditioning
levels.}
\label{fig:size}
\end{figure}

Figure~\ref{fig:size} shows solve times for $n=1000,2000,4000,8000$ and $12000$ at
$p/n=0.5$. The dyadic sizes make the scaling readable directly. Each doubling of $n$
multiplies the time by \DoublingQR\ for QR and \DoublingSchur\ for Schur, against 8 for a
pure cubic, and the fitted log-log slopes are \SlopeQR\ and \SlopeSchur. The operation
count is cubic, but wall time grows slightly slower because the multithreaded BLAS runs well
below peak on $1000\times1000$ matrices and close to it on $12000\times12000$ ones. The
Schur method is faster at every size, and its advantage widens from \SizeRatioLow$\times$
at $n=1000$ to \SizeRatioTop$\times$ at $n=12000$ (\SizeQRTop\ versus \SizeSchurTop~s), consistent
with the constants estimated in Section~\ref{sec:methods}. The seed to seed spread is
small. At $n=12000$ the five QR medians lie between 13.4 and 14.0~s.

The hollow markers repeat the same sizes with $\kap(H)=10^8$ and $\kap(C)=10^6$ and land on
top of the filled ones. Conditioning does not change the number of operations, so it
does not change the time. It changes the accuracy, as the table on the right shows. At
$n=12000$ the QR error grows from $\WellQRXTop$ to $\IllQRXTop$ (four digits lost), the
Schur error from $\WellSchurXTop$ to $\IllSchurXTop$ (ten digits lost), and the Schur
multipliers are off by \IllSchurLTop\%. The gap of seven orders of magnitude is the
same at every size. At $n=1000$ the errors are $\IllQRXLow$ and $\IllSchurXLow$. Because
the reference here is the constructed target, the QR figure is only known to about the
target's own uncertainty, whereas the Schur figure is far above it. The Schur residual
$\eta_{K,F}$, which needs no reference, rises to $\IllSchurKKTTop$ while QR's stays at
$\IllQRKKTTop$.

\subsection{Dependence on the constraint fraction}

\begin{figure}[t]
\centering
\begin{tikzpicture}
\begin{axis}[width=9cm,height=5.6cm,xlabel={Constraint fraction $p/n$},
 ylabel={Complete solve time (s)},xmin=.05,xmax=.95,ymin=0,ymax=\FracYmax,
 xtick={.1,.25,.5,.75,.9},ytick={0,4,8,12,16},legend pos=south east]
\plotqr{v2_fraction_qr.csv}\plotschur{v2_fraction_schur.csv}
\legend{QR/null space,Schur}
\end{axis}\end{tikzpicture}
\caption{Solve time at $n=12000$ against $p/n$, with $\kap(H)=100$ and $\kap(C)=10$, five
seeds and seven batches. Bars show the seed min and max.}
\label{fig:fraction}
\end{figure}

Figure~\ref{fig:fraction} fixes $n=12000$ and varies $p/n$ from $0.1$ to $0.9$. The QR
time is nearly flat, from 12.6 to 15.1~s. Its dominant cost, the full orthogonal factor of
$C^\top$ and the product $HZ$, does not depend much on $p$, and the shrinking reduced system
compensates for the growing basis. The Schur time grows almost linearly from 2.6 to
14.1~s, because the method performs one solve with $H$ per constraint and factors a
$p\times p$ matrix. The speed ratio falls from \FracRatioLow$\times$ at $p/n=0.1$ to
\FracRatioHigh$\times$ at $p/n=0.9$, where the two methods are equal within the seed
spread. No crossover was observed in the tested range, but the ranking is clearly a
function of $p/n$ rather than a property of the methods. Accuracy is unaffected by the
ratio. QR stays at $e_x\approx10^{-14}$ and Schur at $10^{-13}$ throughout.

\subsection{Application problems}

\begin{table}[t]
\centering\footnotesize\setlength{\tabcolsep}{4pt}
\begin{tabular}{@{}lrrcccc@{}}
\toprule
Problem & $n$ & $p$ & QR (s) & Schur (s) & $e_x$ QR & $e_x$ Schur\\
\midrule
DTOC3, original & 12000 & 8000 & \DtocNativeQRT & \DtocNativeSchurT & $\DtocNativeQRX$ & $\DtocNativeSchurX$\\
DTOC3, $x=Uy$ (5 $U$) & 12000 & 8000 & \DtocOrthQRLo\ to \DtocOrthQRHi & \DtocOrthSchurLo\ to \DtocOrthSchurHi & $\DtocOrthQRX$ & $\DtocOrthSchurX$\\
GEMAT1 & 4929 & 5 & \GematFiveQRT & \GematFiveSchurT & $\GematFiveQRX$ & $\GematFiveSchurX$\\
GEMAT1 & 4929 & 20 & \GematTwentyQRT & \GematTwentySchurT & $\GematTwentyQRX$ & $\GematTwentySchurX$\\
Kriging & \KrigNTop & 1 & \KrigQRTop & \KrigSchurTop & \multicolumn{2}{c}{methods agree to $\KrigAgree$}\\
\bottomrule
\end{tabular}
\caption{Application problems. Times are medians over seven batches (and the range over
five orthogonal $U$ for the transformed DTOC3). The DTOC3 reference is a sparse
double precision KKT solve. The GEMAT1 reference is a double precision augmented
least squares solve, itself only accurate to about $10^{-8}$. Kriging has no reference.
Both of its KKT residuals are about $10^{-16}$ and the weights sum to one to $\KrigWsum$.}
\label{tab:apps}
\end{table}

Table~\ref{tab:apps} collects the application problems.

The orthogonally transformed DTOC3 costs the same as the original, 13.0 to 14.0 versus
13.7~s for QR and 8.8 to 9.7 versus 8.7~s for Schur. The solvers do not exploit the diagonal
Hessian, so the DTOC3 timings are a fair dense comparison, and the ranking from the
synthetic problems (Schur about $1.6\times$ faster at $p/n=2/3$) carries over to a
problem with a model behind it. The transformation also preserves the accuracy
picture. QR is three orders of magnitude closer to the reference than Schur in both
representations ($10^{-13}$ versus $10^{-10}$), so the gap is a property of the
algorithms, not of the coordinates.

With very few constraints the Schur method is six to seven times faster, $0.20$ versus
$1.2$~s on GEMAT1 ($p=5$ or $20$) and $1.87$ versus $13.7$~s on kriging ($p=1$). This is the
limit $p\to0$ of Figure~\ref{fig:fraction}. The Schur method reduces to one Cholesky of
$H$, roughly $n^3/3$, while the QR method still forms the full $n\times n$ factor $Q$ and
the product $HZ$ with $Z$ almost square, roughly $2.5n^3$. The predicted ratio of about
$7.5$ matches the measured $7.3$ on kriging.

On GEMAT1 the two methods have the same error, about $10^{-8}$, and I draw no accuracy
conclusion there. The reference is a double precision solve of the original least squares
problem, whereas the solvers receive the explicitly formed $H=A^\top A$ with
$\kap(H)\approx10^9$. The two problems differ by about $\kap(H)\varepsilon\approx10^{-7}$,
and both methods sit at that floor. There is also a structural reason not to expect a
difference. With five well conditioned constraints, $Z$ spans almost all of $\mathbb R^n$,
$G=Z^\top HZ$ has essentially the conditioning of $H$, and the two methods are doing nearly
the same Cholesky.

\subsection{Conditioning}\label{sec:cond}

\begin{figure}[t]
\centering
\begin{tikzpicture}
\begin{axis}[width=7.6cm,height=5.4cm,xmode=log,ymode=log,
 title={Vary $H$, fix $\kap(C)=10$},xlabel={Prescribed $\kap(H)$},ylabel={Relative $x$ error},
 xmin=.5,xmax=3e12,ymin=1e-16,ymax=1e-3,xtick={1,1e4,1e8,1e12},ytick={1e-15,1e-10,1e-5},legend pos=north west]
\plotqr{new_cond_H_qr.csv}\plotschur{new_cond_H_schur.csv}
\legend{QR/null space,Schur}
\end{axis}\end{tikzpicture}
\hfill
\begin{tikzpicture}
\begin{axis}[width=7.6cm,height=5.4cm,xmode=log,ymode=log,
 title={Vary $C$, fix $\kap(H)=100$},xlabel={Prescribed $\kap(C)$},ylabel={Relative $x$ error},
 xmin=.5,xmax=3e12,ymin=1e-16,ymax=1e-3,xtick={1,1e4,1e8,1e12},ytick={1e-15,1e-10,1e-5}]
\plotqr{new_cond_C_qr.csv}\plotschur{new_cond_C_schur.csv}
\node[anchor=north west,text=SchurOrange,font=\footnotesize,align=left]
 at (rel axis cs:.03,.98) {Schur fails 5/5 at $10^{10}$\\and 5/5 at $10^{12}$};
\end{axis}\end{tikzpicture}
\caption{Relative error in $x$ against prescribed condition numbers, $n=80$, $p=20$,
medians and min and max over five seeds, against 96-digit references of the stored problems.
Failed Schur solves have no solution and are not plotted.}
\label{fig:cond}
\end{figure}

Figure~\ref{fig:cond} is the accuracy study with certified references. In the left panel
$H$ becomes ill conditioned while $C$ stays benign. Both errors grow roughly in proportion
to $\kap(H)$, as expected for a backward stable solve, but with different constants. At
$\kap(H)=10^{12}$ the QR median error is $1.1\times10^{-7}$ and the Schur median error
$1.4\times10^{-4}$, and the ratio is between 10 and 100 across the whole sweep. The
multipliers behave differently. Here $e_\lambda$ stays at $10^{-14}$ for QR and
$10^{-12}$ for Schur, because $C$ is well conditioned and $\lambda$ depends on $H$ only
through $g-Hx$.

In the right panel $C$ becomes ill conditioned. Up to $\kap(C)=10^8$ both methods return
solutions, QR again 10 to 100 times more accurate. At $\kap(C)=10^{10}$ and $10^{12}$ the
Schur method fails in all ten cases, because the Cholesky factorization of $S$ encounters
a non positive pivot. The QR method returns a solution in every case, with a median $x$
error of $1.7\times10^{-8}$ at $10^{10}$ and $1.3\times10^{-6}$ at $10^{12}$. It is not
unaffected, however. Its multiplier error at $\kap(C)=10^{10}$ is $92.7\%$ even though
its scaled KKT residual is $1.5\times10^{-17}$. The triangular solve $R\lambda=Q_1^\top(g-Hx)$
inherits the full $\kap(C)$, and a residual at roundoff level says nothing about
$\lambda$ when $R$ is that ill conditioned. A small residual certifies neither component
of the solution on its own.

\subsubsection*{Why the Schur method fails} Equation~\eqref{eq:gram} gives
$\kap(S)=\kap(B)^2$ with $B=CL_H^{-\top}$. At $\kap(C)=10^{10}$ that is of order $10^{20}$,
beyond the $10^{16}$ range of double precision. The smallest eigenvalues of $S$ are then
below the rounding error committed in forming $S=CW$, and the computed $S$ is not
positive definite although the exact one is. To confirm this I repeated the Schur
algorithm, step for step and on the unchanged stored inputs, in \texttt{mpmath} at 64 and
96 decimal digits. All ten failed cases succeed at both precisions, and the solutions
agree with an independent 110-digit solve of the full KKT system to $4\times10^{-80}$. A
double precision reproduction in SciPy fails in all ten cases exactly as MATLAB does, and
the Schur matrices it forms have a negative smallest eigenvalue while differing from the
high precision ones by less than $8\times10^{-16}$ in relative Frobenius norm. The
failures are therefore a finite precision consequence of forming $S$ explicitly, not a
property of the problems. The QR method never forms a product of this kind, and its
orthogonal factor does not square $\kap(C)$.

\subsubsection*{The same picture at $n=12000$} Table~\ref{tab:condlarge} repeats the sweep at
$n=12000$, $p=6000$, with the constructed target as reference. The same ten Schur
failures occur at $\kap(C)\ge10^{10}$, and in the $H$ sweep the gap between the methods is
again four to six digits, with the Schur residual rising to $2\times10^{-6}$ at
$\kap(H)=10^{12}$ while QR's stays at $3\times10^{-17}$. The exact QR digits at large
$\kap$ are limited by the target's own rounding and should not be read as forward
errors. The failures and the residuals need no reference. Times are unchanged at every
condition number, about 14~s for QR and 7~s for Schur.

\begin{table}[t]
\centering\small
\begin{tabular}{@{}lcccccc@{}}
\toprule
 & \multicolumn{3}{c}{Vary $\kap(C)$, $\kap(H)=100$} & \multicolumn{3}{c}{Vary $\kap(H)$, $\kap(C)=10$}\\
\cmidrule(lr){2-4}\cmidrule(lr){5-7}
Target & QR $e_x$ & Schur $e_x$ & Schur ok & QR $e_x$ & Schur $e_x$ & Schur ok\\
\midrule
$10^{4}$ & $2\times10^{-12}$ & $3\times10^{-11}$ & 5/5 & $2\times10^{-13}$ & $1\times10^{-11}$ & 5/5\\
$10^{8}$ & $9\times10^{-9}$ & $1\times10^{-7}$ & 5/5 & $2\times10^{-11}$ & $1\times10^{-7}$ & 5/5\\
$10^{10}$ & $8\times10^{-7}$ & failed & 0/5 & $2\times10^{-10}$ & $1\times10^{-5}$ & 5/5\\
$10^{12}$ & $7\times10^{-5}$ & failed & 0/5 & $2\times10^{-9}$ & $1\times10^{-3}$ & 5/5\\
\bottomrule
\end{tabular}
\caption{Conditioning sweep at $n=12000$, $p=6000$. Medians over five seeds, single solves,
errors against the constructed target.}
\label{tab:condlarge}
\end{table}

\section{Discussion}\label{sec:discussion}

The two methods solve the same system and both cost $O(n^3)$, yet the experiments
separate them cleanly along two axes.

On cost, the Schur complement method wins everywhere in the tested range, and the size of
its advantage is governed by $p/n$, about $2\times$ at $p/n=0.5$, $6$ to $7\times$ when $p$ is
a handful, and nothing at $p/n=0.9$. The reason is visible in the operation counts once
the constants are written out. The Schur method is built from Cholesky factorizations
and triangular solves, the cheapest dense cubic kernels, whereas the QR method must
construct an explicit $n\times n$ orthogonal matrix and multiply $H$ by most of it, and
that cost does not shrink when $p$ does. The dense representation of DTOC3 confirms that
this is not an artefact of structured test matrices.

On accuracy, the QR null space method wins everywhere, by one order of magnitude on
well conditioned problems, by seven orders at $\kap(H)=10^8$, and categorically when $C$
is ill conditioned. The Schur method forms $S=CH^{-1}C^\top$ explicitly and thereby squares
the condition number of $C$, so at $\kap(C)\ge10^{10}$ the computed $S$ is no longer
positive definite in double precision and the method fails outright. The high precision
diagnostic shows that nothing is wrong with the problems or with the algebra. What is
wrong is forming a Gram matrix in floating point. The QR method is not immune to
ill-conditioning in $C$. It returns $x$ to a few digits where Schur returns nothing, but
its multipliers can be meaningless while the KKT residual is at roundoff level, which is
a reminder that a residual, a solution error and a failure count measure different
things and should be reported together.

The practical rule that follows is simple. With few constraints and a well conditioned
$C$, use the Schur complement. With many constraints relative to $n$, or whenever $C$ may
be ill conditioned, use the null space method and pay the factor of two.

Several limitations should be kept in mind. The comparison concerns these two
implementations in MATLAB with MKL, not the methods in the abstract. A Schur method built
on an $LDL^\top$ factorization of $S$, or a null space method using a cheaper basis for
$Z$, would move the numbers. Certified forward errors exist only for the $n=80$ study. At
$n=12000$ the synthetic references are limited by their own construction and the
application references by double precision, so the large-scale accuracy statements are
about residuals, failures and the relative ordering of the two methods rather than exact
digits. The application problems cover a narrow slice of the space (two have very few
constraints, one has a diagonal Hessian in its natural coordinates), and the kriging
covariance parameters were fixed rather than fitted. Finally, all timings come from a
shared cloud environment. The seed to seed spread suggests they are stable, but they
were not repeated across sessions.

\section{Conclusion}

Both dense methods are correct, cubic and, on benign problems, accurate to near machine
precision. The Schur complement method is the faster one by a margin set by $p/n$, and the
QR null space method is the more accurate and the more robust one by a margin set by the
conditioning of $C$. Writing the condition number of the Schur complement as
$\kap(B)^2$ predicts where the Schur method must break in double precision, and a
multiprecision rerun of the failed cases shows that it breaks exactly there and for
exactly that reason.

\section{Code}\label{sec:code}

All solver and experiment code is my own, written in MATLAB, with Python (NumPy, SciPy,
mpmath) used for exporting the test problems, for an independent reimplementation check
of both solvers and for the high precision diagnostic. The complete suite, including the
stored test problems, the raw results of every experiment reported here and the LaTeX
sources, is available at \url{https://github.com/odulsuzkisafilm/qp-dense-solvers}.

\section*{Sources}

The problem formulation and both algorithms follow Benzi, Golub and Liesen~\cite{benzi2005},
Sections 5 to 6. The DTOC3 problem was exported from its SIF definition with
PyCUTEst~1.8.2~\cite{cutest}. The GEMAT1 matrix is from SuiteSparse~\cite{suitesparse} and
the equality constrained least squares setup follows Scott and T\r{u}ma~\cite{scotttuma2022}.
The elevation data are from the UCI repository~\cite{uci_road} (CC~BY~4.0), and the kriging
formulation follows the standard ordinary kriging equations as documented in
GSTools~\cite{gstools}.

\begin{thebibliography}{9}
\bibitem{benzi2005} M.~Benzi, G.~H. Golub and J.~Liesen, Numerical solution of saddle
point problems, \emph{Acta Numerica} 14 (2005) 1 to 137. \href{https://doi.org/10.1017/S0962492904000212}{doi 10.1017/S0962492904000212}.
\bibitem{cutest} N.~I.~M. Gould, D.~Orban and Ph.~L. Toint, CUTEst, a constrained and
unconstrained testing environment with safe threads for mathematical optimization,
\emph{Comput.\ Optim.\ Appl.} 60 (2015) 545 to 557. Problem definitions at
\url{https://github.com/ralna/SIF}.
\bibitem{dtoc3} DTOC3 problem definition, \url{https://github.com/ralna/SIF/blob/master/DTOC3.SIF}.
\bibitem{suitesparse} T.~A. Davis and Y.~Hu, The University of Florida sparse matrix
collection, \emph{ACM Trans.\ Math.\ Softw.} 38 (2011), article 1. Matrix HB/gemat1,
\url{https://sparse.tamu.edu/HB/gemat1}.
\bibitem{scotttuma2022} J.~Scott and M.~T\r{u}ma, Solving large linear least squares
problems with linear equality constraints, \emph{BIT Numer.\ Math.} 62 (2022) 1765 to 1787.
\href{https://doi.org/10.1007/s10543-022-00930-2}{doi 10.1007/s10543-022-00930-2}.
\bibitem{uci_road} M.~Kaul, 3D Road Network (North Jutland, Denmark), UCI Machine Learning
Repository (2013). \href{https://doi.org/10.24432/C5GP51}{doi 10.24432/C5GP51}.
\bibitem{gstools} S.~M\"uller, L.~Sch\"uler et al., GSTools v1.3, a toolbox for geostatistical
modelling in Python, \emph{Geosci.\ Model Dev.} 15 (2022) 3161 to 3182. Ordinary kriging
example in the documentation.
\end{thebibliography}

\end{document}
